\documentclass[10pt, letterpaper]{article}
% ----------------------------------------------------------------------
% LaTeX Preamble
% ----------------------------------------------------------------------

% --- Required Packages ---
\usepackage{reddy_cv}       % Custom CV style package

% ----------------------------------------------------------------------
% Document Body
% ----------------------------------------------------------------------
\begin{document}

% ----------------------------------------------------------------------
% HEADER SECTION
% ----------------------------------------------------------------------
\begin{header}
	\begin{minipage}[t]{1.0\textwidth}
		\fontsize{20pt}{20pt}\selectfont \textbf{Aadarsha Gopala Reddy}\\
		\normalsize
		% Phone
		\mbox{\hrefWithoutArrow{tel:+17408021776}{+1 (740) 802-1776}}%
		\kern \kernSpacing\AND\kern \kernSpacing
		% Email
		\mbox{\hrefWithoutArrow{mailto:adurs2002@gmail.com}{adurs2002@gmail.com}}%
		\kern \kernSpacing\AND\kern \kernSpacing
		\href{https://www.linkedin.com/in/agopalareddy}{LinkedIn}%
		\kern \kernSpacing\AND\kern \kernSpacing
		\href{https://github.com/agopalareddy}{GitHub}%
		\kern \kernSpacing\AND\kern \kernSpacing
		\href{https://agreddy.com}{agreddy.com}%
	\end{minipage}
\end{header}

% --- Vertical Space After Header ---
\vspace{\headerSpacing} % Defined in reddy_cv.sty

% ----------------------------------------------------------------------
% 1. EDUCATION SECTION
% ----------------------------------------------------------------------
\section{EDUCATION}

\begin{educationentry}
	{Washington University in St. Louis} % Institution
	{August 2024 - May 2026} % Date Range
	{M.S. Computer Science} % Degree
	{St. Louis, Missouri, USA} % Location
\begin{highlightsforbulletentries}
\item \textbf{Thesis:} \href{https://openscholarship.wustl.edu/eng_etds/1341/}{Toward Vehicle-Agnostic Driving Signatures for Cognitive Impairment Prediction from Naturalistic Driving Data} (Advisor: Dr. Alvitta Ottley)
\end{highlightsforbulletentries}
\end{educationentry}

\begin{educationentry}
	{Ohio Wesleyan University} % Institution
	{August 2020 - May 2023} % Date Range
	{B.A. Computer Science \& Data Analytics} % Degree
	{Delaware, Ohio, USA} % Location
\begin{highlightsforbulletentries}
\item \textbf{Minor:} Economics
\end{highlightsforbulletentries}
\end{educationentry}

% --- Vertical Space After Section ---
\vspace{\headerSpacing}

% ----------------------------------------------------------------------
% 2. SCHOLARLY WORK SECTION
% ----------------------------------------------------------------------
\section{SCHOLARLY WORK}

\begin{onecolentry}
Gopala Reddy, Aadarsha, ``Toward Vehicle-Agnostic Driving Signatures for Cognitive Impairment Prediction from Naturalistic Driving Data'' (2026). M.S. thesis, Washington University in St. Louis. \href{https://doi.org/10.7936/g95n-fv35}{https://doi.org/10.7936/g95n-fv35}
\end{onecolentry}

% --- Vertical Space After Section ---
\vspace{\headerSpacing}

% ----------------------------------------------------------------------
% 3. RESEARCH EXPERIENCE SECTION
% ----------------------------------------------------------------------
\section{\sectionlink{https://agreddy.com/experience/}{RESEARCH EXPERIENCE}}

\begin{leadershipentry}{Washington University in St. Louis}{St. Louis, Missouri, USA}
	\begin{researchentry}{\href{https://openscholarship.wustl.edu/eng_etds/1341/}{Toward Vehicle-Agnostic Driving Signatures for Cognitive Impairment Prediction from Naturalistic Driving Data}}{August 2025 - May 2026}
		Master's Thesis with the DRIVES Project (PI: Dr. Ganesh Babulal; Advisors: Dr. Alvitta Ottley, Dr. Nathan Jacobs)
		\begin{highlightsforbulletentries}
		\item Built an end-to-end pipeline processing 26,968 participant-weeks of naturalistic driving data from 304 participants, deriving vehicle-agnostic weekly behavioral features and integrating demographic covariates for CDR status prediction.
		\item Compared six CDR prediction models (GRU-DANN, DANN, logistic regression, random forest, XGBoost, MLP) under leave-one-participant-out cross-validation; GRU-DANN achieved highest participant-level ROC AUC (0.599) and balanced accuracy (0.584).
		\end{highlightsforbulletentries}
	\end{researchentry}
\end{leadershipentry}

\begin{leadershipentry}{Ohio Wesleyan University}{Delaware, Ohio, USA}
	\begin{researchentry}{\href{https://github.com/agopalareddy/DATA490}{Opinion Survey on Artificial Intelligence in the Workplace}}{January 2023 - May 2023}
		Independent Study (Mentor: Dr. Nicholas Dietrich)
		\begin{highlightsforbulletentries}
		\item Designed, fielded, and investigated a survey experiment on the impact of AI on 250 employees within each of four industries using survey data.
		\item Analyzed data using Tableau dashboards and R to identify trends and common opinions.
		\item Presented preliminary results on AI's positive impact on employee work at the OWU Spring Student Symposium.
		\end{highlightsforbulletentries}
	\end{researchentry}
	\begin{researchentry}{\href{https://agreddy.com/lost-cities/}{Artificial Intelligence in Modern Board Games --- Lost Cities AI}}{May 2022 - July 2022}
		Summer Science Research Program (Mentor: Dr. Sean McCulloch)
		\begin{highlightsforbulletentries}
		\item Developed a digital version of the Lost Cities card game using Java.
		\item Implemented an intelligent agent that played the game against a human, winning 13 of 18 games in testing.
		\end{highlightsforbulletentries}
	\end{researchentry}
\end{leadershipentry}

% --- Vertical Space After Section ---
\vspace{\headerSpacing}

% ----------------------------------------------------------------------
% 4. TEACHING \& MENTORSHIP EXPERIENCE SECTION
% ----------------------------------------------------------------------
\needspace{16\baselineskip}
\section{\sectionlink{https://agreddy.com/experience/}{TEACHING \& MENTORSHIP EXPERIENCE}}

\begin{leadershipentry}{Washington University in St. Louis}{St. Louis, Missouri, USA}
	\begin{positionentry}{Graduate Teaching Assistant - CSE 5114: Data Manipulation and Management at Scale}{January 2026 - May 2026}
		\item Taught large-scale data pipeline concepts (Snowflake, Airflow, Spark, Kafka, Flink) to a graduate cohort, covering fault-tolerant design, streaming windowing strategies, and query optimization.
		\item Designed oral midterms as mock technical interviews, requiring students to reason through pipeline architecture decisions and troubleshoot system failures under real-world constraints.
		\item Mentored three capstone teams from raw open datasets to production-grade dashboards, providing weekly feedback that kept all teams on schedule through final presentations.
		\item Held weekly office hours to debug student pipeline implementations across Snowflake schema design, Airflow DAG failures, and Kafka consumer configuration.
	\end{positionentry}
	\needspace{6\baselineskip}
	\begin{positionentry}{HackWashU Databases Workshop}{October 2025}
		\item Designed and facilitated a hands-on databases workshop for beginner-to-intermediate developers, teaching through two complete projects rather than lectures.
		\item Guided participants through building a normalized SQLite database from the iTunes Search API and a full-stack Supabase Todo app with Row Level Security — covering relational databases, APIs, cloud Postgres, and access control.
	\end{positionentry}
	\needspace{4\baselineskip}
	\begin{positionentry}{Engineering Workshop by AGES \& Science Coach}{September 2024}
		\item Delivered an introductory workshop on artificial intelligence to high school students, covering core AI concepts with hands-on demonstrations.
	\end{positionentry}
\end{leadershipentry}

\needspace{8\baselineskip}
\begin{leadershipentry}{Ohio Wesleyan University}{Delaware, Ohio, USA}
	\begin{positionentry}{Computer Science Lab Assistant}{January 2022 - May 2023}
		\item Tutored approximately ten students for over 6 hours per week, assisting with homework and exam preparation for introductory, intermediate, and advanced computer science coursework.
		\item Tested and graded assignments of around 35 students in introductory and intermediate computer science classes.
	\end{positionentry}
\end{leadershipentry}

% --- Vertical Space After Section ---
\vspace{\headerSpacing}

% ----------------------------------------------------------------------
% 5. PRESENTATIONS SECTION
% ----------------------------------------------------------------------
\needspace{5\baselineskip}
\section{PRESENTATIONS}

\begin{leadershipentry}{Ohio Wesleyan University}{Delaware, Ohio, USA}
	\begin{presentationentry}{Spring Student Symposium}{April 2023}
		Artificial Intelligence Opinion Survey
	\end{presentationentry}
	\begin{presentationentry}{Patricia Belt Conrades Summer Science Research Symposium}{September 2022}
		Artificial Intelligence in Modern Board Games
	\end{presentationentry}
\end{leadershipentry}

% --- Vertical Space After Section ---
\vspace{\headerSpacing}

% ----------------------------------------------------------------------
% 6. GRANTS, FELLOWSHIPS, \& AWARDS SECTION
% ----------------------------------------------------------------------
\section{GRANTS, FELLOWSHIPS, \& AWARDS}

\begin{honorentry}
	{April 2026} % Date/Year
	{\textbf{Bauer Leaders Academy (BLA) Leadership Development Badge}, Washington University in St. Louis} % Honor Name & Institution
\end{honorentry}

\begin{honorentry}
	{May 2024} % Date/Year
	{\textbf{Affiliated School Scholarship}, Washington University in St. Louis} % Honor Name & Institution
\end{honorentry}

\begin{honorentry}
	{May 2023} % Date/Year
	{\textbf{Dean's List}, Ohio Wesleyan University} % Honor Name & Institution
\end{honorentry}

\begin{honorentry}
	{April 2023} % Date/Year
	{\textbf{Mortar Board National College Senior Honor Society Membership}, Ohio Wesleyan University} % Honor Name & Institution
\end{honorentry}

\begin{honorentry}
	{April 2023} % Date/Year
	{\textbf{Golden Bishop Award - WCSA Best New Member}, Ohio Wesleyan University} % Honor Name & Institution
\end{honorentry}

\begin{honorentry}
	{May 2022} % Date/Year
	{\textbf{Dean's List}, Ohio Wesleyan University} % Honor Name & Institution
\end{honorentry}

\begin{honorentry}
	{April 2022} % Date/Year
	{\textbf{Florence Leas Sophomore Prize}, Ohio Wesleyan University} % Honor Name & Institution
\end{honorentry}

\begin{honorentry}
	{March 2022} % Date/Year
	{\textbf{First Place - 32nd Annual Spring Programming Contest}, Denison University} % Honor Name & Institution
\end{honorentry}

\begin{honorentry}
	{January 2020} % Date/Year
	{\textbf{International Baccalaureate Scholarship}, Ohio Wesleyan University} % Honor Name & Institution
\end{honorentry}

\begin{honorentry}
	{January 2020} % Date/Year
	{\textbf{Scholar}, Next Genius Scholarship Foundation} % Honor Name & Institution
\end{honorentry}

% --- Vertical Space After Section ---
\vspace{\headerSpacing}

% ----------------------------------------------------------------------
% 7. PROFESSIONAL SERVICE \& LEADERSHIP SECTION
% ----------------------------------------------------------------------
\section{PROFESSIONAL SERVICE \& LEADERSHIP}

\begin{leadershipentry}{Washington University in St. Louis}{St. Louis, Missouri, USA}
	\begin{positionentry}{Vice President of the Graduate-Professional Council (GPC) Chamber, Graduate and Professional Student Council (GPSC)}{May 2025 - May 2026}
		\item Built and maintained committee membership rosters using ranked-choice preference allocation to balance workload and ensure efficient cross-committee coordination.
		\item Oversaw information architecture for the Council's documentation repository.
		\item Represented the graduate student body to university leadership and supported the GPC President in advancing strategic council priorities.
		\item Chaired the Constitution Committee (August 2025 - February 2026), coordinating stakeholder input to draft and ratify GPSC governing documents aligned with the Council's mission and university policies.
	\end{positionentry}
	\needspace{8\baselineskip}
	\begin{positionentry}{Member, Graduate Student Affairs Advisory Board (GSAAB)}{August 2025 - May 2026}
		\item Provided strategic feedback to the Vice Chancellor for Student Affairs on critical issues impacting graduate student life and experience.
		\item Engaged in critical dialogue with senior university leaders, including the Dean of Students, to advocate for improvements to the graduate student experience across all academic programs.
		\item Contributed to monthly discussions on student affairs policies and initiatives to enhance support systems and resources for the graduate student community.
	\end{positionentry}
	\needspace{8\baselineskip}
	\begin{positionentry}{Graduate Student Member, Center for Career Engagement (CCE) Student Advisory Board}{August 2025 - May 2026}
		\item Guided the strategic direction of the Center for Career Engagement to ensure services align with diverse needs of graduate students from all academic backgrounds and career interests.
		\item Provided insights on career engagement programs, workshops, and resources to enhance accessibility and effectiveness for the WashU graduate student community.
		\item Built community around career development initiatives and fostered collaboration between the CCE and graduate student body.
	\end{positionentry}
	\needspace{3\baselineskip}
	\begin{positionentry}{Treasurer, Umang (Indian Graduate Student Association)}{June 2025 - October 2025}
		\item Managed club finances and organized cultural events and festivals, promoting cultural awareness and fostering community among Indian graduate students.
	\end{positionentry}
	\begin{positionentry}{Treasurer, HackWashU}{December 2024 - July 2025}
		\item Managed the club's finances, including budgeting and expense tracking, to ensure sustainable operations and successful event planning.
		\item Coordinated fundraising efforts and sponsorship outreach to support club activities and initiatives.
		\item Collaborated with club members to organize hackathons and coding events, fostering a culture of innovation and teamwork among participants.
	\end{positionentry}
\end{leadershipentry}

\begin{leadershipentry}{Ohio Wesleyan University}{Delaware, Ohio, USA}
	\begin{positionentry}{Budget Committee Senator, Wesleyan Council on Student Affairs (WCSA)}{August 2022 - May 2023}
		\item Proposed and voted on senate bills while overseeing allocation of an annual budget exceeding \$350,000 to campus individuals and organizations.
		\item Received the Golden Bishop Award for WCSA Best New Member for overall contribution to WCSA's mission and service on the Budget Committee.
	\end{positionentry}
	\begin{positionentry}{Member, The Neurds (Neuroscience Club)}{August 2022 - May 2023}
		\item Presented on AI at Brain Fair 2023 during Brain Awareness Week, bridging neuroscience and technology.
	\end{positionentry}
	\begin{positionentry}{Member, Mathematics, Computer Science \& Data Analytics Student Board}{August 2021 - May 2023}
		\item Analyzed faculty evaluations and synthesized student feedback on course offerings to recommend departmental improvements.
	\end{positionentry}
	\begin{positionentry}{Vice President of Finance, Campus Programming Board}{January 2022 - December 2022}
		\item Managed annual budgets exceeding \$86,000 for campus-wide programming and secured \$10,000 in additional annual funding through negotiation with WCSA.
		\item Revised the organization's governing documents, codifying financial and operational practices.
	\end{positionentry}
\end{leadershipentry}

% --- Vertical Space After Section ---
\vspace{\headerSpacing}

% ----------------------------------------------------------------------
% 8. PROJECTS SECTION
% ----------------------------------------------------------------------
\needspace{10\baselineskip}
\section{\sectionlink{https://agreddy.com/projects/}{PROJECTS}}

\begin{leadershipentry}{Washington University in St. Louis}{St. Louis, Missouri, USA}
	\begin{positionentry}{\href{https://github.com/agopalareddy/CSE5519-Project}{Red-Blue Visual Auto Defender: Automated Visual Jailbreak Generation and Explainable Defenses}}{September 2025 - October 2025}
		\item Built an automated red-blue teaming pipeline for vision-language model (VLM) security with teammates Stuart Aldrich and Mohammad Rouie Miab, addressing visual prompt injection (VPI) attacks against VLM-based agents.
		\item Generated visual jailbreak attack images by overlaying malicious instructions onto benign images and evaluated them against a target VLM (Gemma-3-4b-it) using semantic response analysis.
		\item Auto-generated deterministic, OCR-based Python defenses for each successful attack, validated with accuracy, precision, and recall metrics in an iterative refinement loop.
	\end{positionentry}
	\begin{positionentry}{\href{https://github.com/agopalareddy/CSE419A_Multimodal_Prediction_of_Alzheimers}{Multimodal Prediction of Alzheimer's Disease}}{August 2024 - December 2024}
		\item Developed a multimodal fusion approach for early detection of Alzheimer's Disease, combining MRI imaging and clinical/demographic data from the OASIS-1 dataset.
		\item Implemented CNNs using TensorFlow/Keras for brain MRI analysis alongside XGBoost and Random Forest models for tabular clinical data.
		\item Created a combined classifier integrating both modalities and documented the project in NeurIPS format with comprehensive performance evaluation metrics and visualizations.
	\end{positionentry}
	\begin{positionentry}{\href{https://github.com/agopalareddy/CSE510A_Datacenter_Cooling}{Datacenter Cooling Optimization using Deep Reinforcement Learning}}{August 2024 - December 2024}
		\item Implemented multiple deep reinforcement learning algorithms (DDQN, PPO, SAC) integrated with EnergyPlus simulations to optimize datacenter cooling systems.
		\item Achieved 35.8\% energy efficiency improvement over baseline using DDQN, addressing the underserved small to mid-sized datacenter market.
		\item Created a comprehensive framework integrating building energy simulation with reinforcement learning for dynamic thermal management.
	\end{positionentry}
\end{leadershipentry}

\needspace{8\baselineskip}
\begin{leadershipentry}{MITxSureStart FutureMakers Create-a-Thon Program}{Remote}
	\begin{positionentry}{\href{https://github.com/agopalareddy/AIPS}{Parkinson's Disease AI Diagnosis Software}}{June 2022 - July 2022}
		\item Built a dual-input pipeline combining audio spectrograms (custom CNN for vocal tremor detection) and spiral-drawing images (Random Forest classifier with augmentation for micrographia detection).
		\item Implemented a weighted ensemble combining both modalities into a single Parkinson's diagnosis score.
		\item Designed a Flask web interface for user submissions of audio and drawing samples for real-time risk assessment.
	\end{positionentry}
\end{leadershipentry}

% --- Vertical Space After Section ---
\vspace{\headerSpacing}

% ----------------------------------------------------------------------
% 9. PROFESSIONAL EXPERIENCE SECTION
% ----------------------------------------------------------------------
\section{\sectionlink{https://agreddy.com/experience/}{OTHER PROFESSIONAL EXPERIENCE}}

\begin{leadershipentry}{Crittero, Inc.}{Remote}
	\begin{positionentry}{AI Engineer Intern}{June 2025 - August 2025}
		\item Generated labeled social media training data for seven user personas via time-series behavioral modeling in TypeScript, capturing post sequences, engagement timings, and content preferences.
		\item Trained a TensorFlow/Keras recommendation model with persona-type, content-category, and temporal embeddings, plus a collaborative-filtering cold-start fallback for new users.
		\item Built a configurable Python CLI for reproducible model evaluation, enabling systematic comparison across hyperparameter settings, embedding dimensions, and fallback thresholds.
	\end{positionentry}
\end{leadershipentry}

\begin{leadershipentry}{Washington University in St. Louis}{St. Louis, Missouri, USA}
	\begin{positionentry}{Graduate Assistant, Taylor Family Center for Student Success}{August 2024 - May 2026}
		\item Planned and ran three recurring engagement programs for first-generation, limited-income (FLI) students — Welcome Dinner, First-Gen Week, and End of Week Unwind — coordinating logistics, promotion, and on-the-day operations.
		\item Authored FERPA-compliant process manuals for engagement data collection and evaluation, giving center staff a standardized, repeatable framework for measuring FLI student participation.
		\item Analyzed semester survey responses and event attendance data, producing reports that directly informed the following cycle's programming decisions and outreach priorities.
		\item Managed the center's weekly newsletter; represented the department at off-campus excursions and FLI student trips.
	\end{positionentry}
\end{leadershipentry}

\begin{leadershipentry}{Lab714}{Boca Raton, Florida, USA}
	\begin{positionentry}{Data Analyst Intern}{June 2023 - May 2024}
		\item Analyzed IoT environmental sensor data streams on AWS, building dashboards that surfaced chemical imbalances and equipment anomalies through automated threshold detection and trend analysis.
		\item Built AWS data ingestion pipelines to normalize and store continuous sensor telemetry from deployed devices, supporting real-time operational monitoring of deployed IoT fleets.
		\item Wrote SQL queries to aggregate historical sensor data by pool, date range, and chemical parameter, powering the trend analytics and anomaly detection shown in client dashboards.
		\item Designed and engineered the Figma-to-React UI for an all-in-one environmental sensor device, delivering a pixel-perfect, production-ready interface from initial wireframes.
	\end{positionentry}
\end{leadershipentry}

% --- Vertical Space After Section ---
\vspace{\headerSpacing}

% ----------------------------------------------------------------------
% 10. TECHNICAL SKILLS \& AFFILIATIONS SECTION
% ----------------------------------------------------------------------
\section{TECHNICAL SKILLS \& LANGUAGES}

\begin{skillcategory}{Programming Languages}
	Java, C++, Python, R, JavaScript/TypeScript, PHP, C\#, SQL, Rust.
\end{skillcategory}

\begin{skillcategory}{Frameworks \& Libraries}
	TensorFlow, PyTorch, Scikit-learn, Keras, NumPy, Pandas, Matplotlib, Seaborn, Node.js, Express.js, Vue.js, React.
\end{skillcategory}

\begin{skillcategory}{Tools \& Technologies}
	\LaTeX, Git, Tableau, Power BI, MySQL, AWS, MongoDB, Snowflake, Apache Airflow, Spark, Kafka, Flink.
\end{skillcategory}

\begin{skillcategory}{Languages}
	Proficient in English, Kannada, Telugu; Conversational in Hindi.
\end{skillcategory}

% ----------------------------------------------------------------------
% End of Document
% ----------------------------------------------------------------------
\end{document}
